\section*{Abstract}
A 2D foam coarsens over time, larger bubbles grow and smaller bubbles disappear. As this happens the foam restructures and the bubbles get displaced. This paper investigates the statistical properties of individual bubbles using tools borrowed from econophysics.  Drawing a basis between bubble displacements and financial price movements, we apply log returns, volatility, temporal autocorrelation, complementary cumulative distribution functions (CCDFs), and mean squared displacement (MSD). Throughout this investigation, Brownian motion is used as a baseline, comparing with this pure randomness to characterise the particular type of randomness of these bubbles movements. We find that it is hard to draw conclusions from the bubble MSD so move on to more statistical methods to analyse. Bubble displacement distributions have a sharp peak around zero and taper off each side however for large $\Delta t$, looking at dispalcement over 10, 50, 100 timesteps or so, the distribution approaches a Gaussian curve. These distributions do not scale accoring to a Levy distribution like the financial distributions. The CCDF cut depends on where the cut is taken rather than being linear. A two parameter function of the form $f(C, \Delta t) = c_1(C) \Delta t/ (1 + c_2(C)\Delta t)$ is proposed to describe the CCDF scaling. Notably, the autocorrelation function (ACF) of the log returns remains positive for the whole simulation, while the volatility ACF follows a power law decay similar to that of the financial data. These findings suggest that the bubble dynamics in a coarsening foam follow a particular non trivial statistical stucture beyond simple Brownian motion.